% ============================================
% Implementing FRFP in Scientific Research
% ============================================
\chapter{Case Studies: Implementing FRFP in Real Systems}
\label{ch:frfp-implementations}
This chapter translates the FRFP core into concrete implementation patterns across four domains:
scientific research, engineering workflows, financial decision-making, and clinical practice. The aim
is practical rather than theoretical: given a real workflow, identify the explicit artefacts the AI
can touch, locate the boundary events where those artefacts become endorsed commitments, and design
safe horizons and governance so that long-horizon use does not quietly convert convenience into
unaccountable automation. The sections that follow are not exhaustive “best practices” lists. They
are domain-specific instantiations of the same derivation logic: make commitment points explicit,
make the evidence bundle at each boundary visible and loggable, impose finite limits that force
periodic re-grounding, and specify who has authority to tighten horizons, halt modes, and revise the
protocol after incidents.
\section{How to Derive FRFP Implementation Guidelines in Any Domain}
\label{sec:frfp-derive-guidelines}

This section is not a checklist of best practices. It is a method for turning the FRFP core objects
into domain-specific operational controls. The derivation procedure is the same in every domain:
start from the explicit/tacit split, treat the real system as runs that accumulate over time, and
design boundaries, safe horizons, and governance so that endorsement remains meaningfully human and
auditable.

\subsection{Derivation Procedure}

To produce an FRFP implementation for a new domain, follow these steps.

\begin{itemize}
  \item \textbf{D1: Fix the unit of commitment.}
        Identify the domain’s real commitment points: the moments where an explicit artefact becomes a
        binding action, claim, or allocation of resources. Typical examples include: placing an order,
        signing a report, filing a document, approving a release, committing capital, or publishing a
        scientific claim. These are your boundary candidates.

  \item \textbf{D2: Enumerate the explicit artefact types.}
        Define what counts as an explicit artefact in the domain: inputs, intermediate objects, model
        outputs, dashboards, reports, drafts, logs, and any structured representations used for decision.
        Separate artefacts that are \emph{descriptive} (summaries, extractions) from artefacts that are
        \emph{decision-shaping} (rankings, risk scores, recommended actions), since the latter require
        stricter boundary controls.

  \item \textbf{D3: Name the tacit evaluators and their standards.}
        Specify who bears correctness and responsibility at each commitment point and what tacit
        standards they use (professional norms, legal duties, safety criteria, epistemic standards).
        Include authority, time budget, and incentive conflicts. If meaningful evaluation is infeasible
        in the real workflow, the protocol is misdesigned regardless of model quality.

  \item \textbf{D4: Trace runs, not snapshots.}
        Draw representative run families: normal operation, degraded operation (missing data, drift,
        outages), adversarial operation (gaming, manipulation), and growth runs (scaling volume, new
        populations, new settings). Mark where history matters: feedback loops, accumulation of exposure,
        habituation, and institutional adaptation.

  \item \textbf{D5: Engineer boundaries as explicit events.}
        For each commitment point, specify the boundary event as an explicit protocol step with:
        (i) required artefacts shown to the endorser, (ii) required provenance (what generated what),
        (iii) required checks (consistency, contraindications, guardrails), (iv) a logged decision record
        (accept/reject/modify + rationale category), and (v) escalation and override rules. A boundary is
        not ``a human is in the loop''; it is a designed event that produces reconstructible evidence.

  \item \textbf{D6: Define safe horizons as enforceable constraints.}
        Convert long-horizon risk into enforceable limits: caps on consecutive uses, caps on binding
        actions per configuration, periodic re-validation requirements, and drift triggers that force
        downgrade to advisory-only mode. State who can trigger resets, what evidence is required, and
        what changes (halt, rollback, audit, re-training, re-approval). Horizons should be tighter for
        high-stakes or hard-to-detect failure modes.

  \item \textbf{D7: Specify governance operators and change-control.}
        Name the bodies that can change the system: who can modify thresholds, prompts, UI affordances,
        model versions, and workflow rules; who can tighten horizons; who can suspend operation; and how
        incidents force protocol revision. Require at least one independent channel of evaluation
        (audit, expert review, external assessment, or regulator access) that is not reducible to the
        system’s own dashboards.

  \item \textbf{D8: Pre-register failure playbooks.}
        For each major risk indicator, pre-commit the response: what triggers a freeze, downgrade,
        mandatory audit, rollback, or halt; who has authority; and what must be documented. This prevents
        improvisation under pressure and makes safe horizons enforceable.

  \item \textbf{D9: Produce the domain implementation outputs.}
        Every FRFP implementation should yield concrete outputs:
        (i) a typed artefact inventory (what the AI may produce and what it may not),
        (ii) a boundary specification (events, evidence bundles, logs),
        (iii) safe-horizon policies (caps, triggers, resets),
        (iv) a run-aware monitoring profile (signals of drift and boundary hollowing-out),
        and (v) governance and change-control rules (who can do what, when, and with what evidence).
\end{itemize}

\subsection{How to Use This Section}

A reader from a new domain can apply the procedure above in two passes. First, produce a minimal
implementation by identifying one run family and the top two commitment boundaries, then define the
logs and horizon triggers required to keep endorsement meaningful. Second, expand coverage to degraded,
adversarial, and growth runs, tightening horizons and governance where the domain’s stakes and failure
modes demand it. The implementation patterns in the rest of this chapter are concrete instantiations
of this derivation method in science, engineering, finance, and clinical practice.

\section{Implementing FRFP in Scientific Research Workflows}
\label{sec:frfp-science-implementation}

This section develops a concrete implementation of the Foundational Reasoning and Feedback Protocol
(FRFP) for scientific research. The core idea is simple: AI systems can generate and transform explicit
research artefacts, but correctness and endorsement remain human responsibilities. FRFP therefore treats
scientific work as long-horizon collaboration where the critical design objects are (i) the explicit
workspace of artefacts and pipelines, (ii) the tacit judgments that determine what counts as accepted
knowledge, and (iii) the boundary events where outputs become endorsed claims, results, or decisions.
We specialise these ideas to day-to-day research and describe practical instrumentation for runs that
unfold over weeks or months.

\subsection{Scope and Objectives}

We focus on research settings where:

\begin{itemize}
  \item large language models and tool-using systems assist with literature
        review, ideation, proof sketches, modelling, simulation, and drafting;
  \item correctness still ultimately matters in a strong sense:
        scientific claims, theorems, and experimental conclusions are expected
        to withstand scrutiny and serve as a basis for further work;
  \item collaboration is long-horizon: projects extend over weeks or months,
        with many intermediate artefacts and partial results.
\end{itemize}

An FRFP-aligned implementation for scientific research should:

\begin{enumerate}
\item respect the explicit/tacit separation:
        AI systems operate only on explicit artefacts; correctness authority resides in human judgment
        at clear endorsement points;
  \item make long-horizon pipelines auditable (what was produced, what was checked, what was endorsed);
  \item manage degradation of tacit grounding through safe horizons and required re-grounding;
  \item provide instruments (logs, ledgers, metrics) that researchers and committees can actually use.

\end{enumerate}

\subsection{Core Design Principles for Scientific Workflows}

We state core design principles as implementation constraints on FRFP runs
specific to scientific research.

\subsubsection*{P1: Explicit Research Workspace}

All research artefacts that can, in principle, be made explicit are treated as
elements of explicit space $E$:

\begin{itemize}
  \item literature summaries, citation graphs, and topic maps;
  \item conjectures, proof sketches, formal definitions, models, and code;
  \item experiment designs, data processing pipelines, and results;
  \item manuscripts, referee reports, and internal notes.
\end{itemize}

AI systems are allowed to read, generate, transform, and organise these
artefacts.
Their actions are explicit transformations over the workspace: drafting, summarising, refactoring,
reformatting, and proposing candidate structures.

\subsubsection*{P2: Tacit Scientific Understanding and Endorsement}

Tacit space $T$ contains the epistemic states of human researchers:

\begin{itemize}
  \item their background knowledge, intuitions, and sense of plausibility;
  \item their understanding of methods, limitations, and domain-specific norms;
  \item their individual and collective judgments about whether a result
        is acceptable as a basis for further work.
\end{itemize}

FRFP enforces that correctness of scientific claims is attached only after
human evaluation:
accepting a theorem, a result, or a modelling choice is always a human act of
endorsement at a boundary.

\subsubsection*{P3: Explicit Boundary Events in Research}

For a scientific project, boundary events include:

\begin{itemize}
  \item accepting a conjecture as sufficiently precise to pursue seriously;
  \item signing off on a proof sketch as provisionally correct;
  \item freezing a model or experiment design as the basis for data collection;
  \item submitting a manuscript or internal report;
  \item endorsing a result as part of a review, panel, or committee decision.
\end{itemize}

An FRFP implementation:

\begin{itemize}
  \item designates specific points in the research pipeline as boundary events;
  \item identifies the responsible role(s) (e.g.\ lead author, PI, committee);
  \item records the explicit artefacts and context available at each boundary (inputs, outputs,
        provenance pointers, and any checks performed);
  \item treats these events as the points where correctness in $T$ is attached.
\end{itemize}

\subsubsection*{P4: Safe Horizons and Re-Grounding}

Scientific research is iterative and extended.
Within FRFP, each project is associated with \emph{safe horizons}:

\begin{itemize}
  \item limits on how many AI-assisted transformations can occur without
        returning to primary material (original papers, raw data, base axioms);
  \item limits on how many endorsement steps can be chained without independent
        peer review or re-derivation;
  \item constraints on how far formal manipulations can drift from motivating
        questions without an explicit re-alignment step.
\end{itemize}

Implementation must specify:

\begin{itemize}
  \item when researchers must read primary sources rather than summaries;
  \item when proofs or derivations must be reconstructed without AI assistance;
  \item when external or independent review is required.
\end{itemize}

These safe-horizon constraints are treated as part of the admissibility
conditions on runs.

\subsection{FRFP Research Methodology as Pipelines}
\label{sec:frfp-methodology-research}

We now describe a concrete research methodology as patterns of explicit
pipelines and boundary events.

\subsubsection*{M1: Context Management in Research Pipelines}

For a given project, treat the evolving work as a run: a sequence of research actions that create,
transform, and promote explicit artefacts over time. The point is not the notation; it is the discipline
that the project has a traceable history of what depended on what, and where endorsement occurred.


Each non-trivial action should carry a context annotation:

\begin{itemize}
  \item identifiers of upstream artefacts on which it depends
        (e.g.\ which papers, lemmas, or datasets);
  \item the version of key definitions and notation in force at that point;
  \item the project intent and phase (e.g.\ exploratory, formalisation,
        validation).
\end{itemize}

\paragraph{Implementation guidelines.}

\begin{itemize}
\item Maintain a versioned project graph in which nodes are artefacts and edges are transformations;
        ensure each non-trivial step records its parent artefacts and context.
  \item At boundary events, present human evaluators not only with the target
        artefact but with the relevant slice of this graph, so that tacit
        judgment can take context into account.
  \item Provide tools to ``slice'' the run by section, claim, or experiment,
        so that localized audits of correctness and provenance are feasible.
\end{itemize}

\subsubsection*{M2: Iterative Co-Formalization of Scientific Questions}

Scientific projects often begin from loosely specified questions and gradually
move toward formal statements.
Within FRFP, this is a navigation pattern across levels of formality:

\begin{itemize}
\item moves between viewpoints and granularity (informal description, proto-formal statement,
        full formalisation);
  \item explicit edits refine definitions, introduce notation, and connect informal claims to checkable
        objects (proof obligations, data checks, derivations, or cited evidence).
\end{itemize}

\paragraph{Phase labels.}
It is useful to tag each segment of the project run with a \emph{formality phase}:

\begin{itemize}
  \item F0: informal exploration (questions, intuitions, examples);
  \item F1: proto-formal (candidate definitions, rough conjectures);
  \item F2: formal development (definitions, lemmas, proofs, models);
  \item F3: consolidation (cleaned-up proofs, canonical statements,
        publication-ready artefacts).
\end{itemize}

\paragraph{Implementation guidelines.}

\begin{itemize}
\item Require an explicit recorded transition for phase changes
      (e.g.\ ``promote this informal claim to a conjecture with this statement'').
  \item At higher phases (F2/F3), restrict admissible AI operations:
        for example, proof steps may be suggested but require local checking,
        and model refactorings must preserve stated invariants.
  \item Ensure that there are navigation links back from each formal artefact
        to its informal motivation, supporting tacit understanding at the
        boundary.
\end{itemize}

\subsubsection*{M3: Reproducible Human--AI Research Sessions}

Extended co-reasoning with AI systems---for example, over many days of proof
search or model refinement---should be treated as structured runs, not as
ephemeral chats.

\paragraph{Session structure.}
We partition a project into \emph{sessions} $(S_1, S_2, \dots)$.
Each session $S_k$ corresponds to a run segment: a bounded interval of work with its own intent,
artefacts, and endorsement points.

\begin{itemize}
  \item an explicit session intent (e.g.\ ``clarify definitions in Section~2'',
        ``search for counterexamples to Lemma~3.4'');
  \item a bounded safe horizon (e.g.\ number of AI-assisted steps or time);
  \item one or more boundary events (e.g.\ acceptance of a refined statement,
        rejection of a proposed lemma, escalation of an open problem).
\end{itemize}

\paragraph{Implementation guidelines.}

\begin{itemize}
  \item Log each session as a coherent run segment with:
        prompts, visible AI outputs, human edits, navigation actions,
        and boundary decisions.
  \item At the end of a session, write a short human-authored summary that
        interprets the explicit progress in tacit terms:
        what was learned, what remains uncertain, what changed in the
        researcher's understanding.
  \item Support replay at two levels:
        raw interaction (for debugging and training) and
        structural pipeline view (for scientific auditing and reuse).
\end{itemize}

\subsection{Instruments for FRFP-Aligned Scientific Practice}
\label{sec:frfp-instruments-science}

We now introduce a small set of instruments specialised to research settings.

\subsubsection*{I1: Research Integrity Ledger}

The \emph{research integrity ledger} is a structured record of a project's
status relative to FRFP constraints.

For each significant pipeline segment (e.g.\ a section, a cluster of lemmas,
an experiment series), the ledger stores:

\begin{itemize}
  \item phase labels (F0--F3) and project intent;
  \item an integrity triple $(V, C, P)$:
        \begin{itemize}
          \item $V$: structural validity of the pipeline (protocol compliance,
                type discipline, absence of gross methodological violations);
          \item $C$: tacit correctness status as judged by responsible
                experts at boundary points (e.g.\ ``provisionally accepted'',
                ``rejected'', ``open'');
          \item $P$: provenance adequacy (traceable sources, clear use of
                AI assistance, accessible upstream artefacts).
        \end{itemize}
  \item links to boundary events (who endorsed what, when, under which
        conditions).
\end{itemize}

\paragraph{Implementation.}

\begin{itemize}
  \item Index ledger entries by project, section, and artefact ID.
  \item Provide views that answer questions such as:
        \begin{itemize}
          \item ``Which results in this paper are structurally valid but still
                open in $C$?''
          \item ``Which parts of the proof relied heavily on AI suggestions?''
          \item ``Where is provenance currently weak or incomplete?''
        \end{itemize}
  \item Treat ledger updates as first-class events: updating $C$ or $P$
        in response to new evidence is itself a logged action.
\end{itemize}

\subsubsection*{I2: Hybrid Authorship Log}

The \emph{hybrid authorship log} records, at a fine-grained level, how AI and
human contributions are interleaved in research artefacts.

For each document, proof, or model, the log stores:

\begin{itemize}
  \item a sequence of edits (diffs) with author type labels
        (human, AI, tool-assisted, mixed);
  \item pointers from each region of the artefact (e.g.\ a paragraph, lemma,
        or code block) to the chain of edits that produced it;
  \item links from edits back to the corresponding steps in the FRFP run.
\end{itemize}

\paragraph{Implementation.}

\begin{itemize}
  \item Integrate authorship logging into the tools through which AI assistance
        is provided (e.g.\ IDEs, notebooks, document editors).
  \item Expose summarised authorship information at boundary events:
        when a researcher endorses a theorem or draft, they can see which parts
        are predominantly AI-authored and adjust their scrutiny accordingly.
  \item Support post-hoc analysis of projects: for example, computing statistics
        on where AI contributions cluster (definitions vs.\ proofs vs.\ exposition).
\end{itemize}

\subsubsection*{I3: Peer and Committee Evaluation Logs}

Scientific correctness is often mediated by peer review and committee decisions.
Within FRFP, these are institutional boundary events combining many tacit
evaluations.

For each review or committee decision, an evaluation log records:

\begin{itemize}
  \item the explicit artefacts considered (manuscript versions, proof logs,
        datasets, models);
  \item individual judgments (accept / reject / request changes; confidence);
  \item the institutional rule used to combine these judgments (simple majority,
        weighted vote, veto, chair decision);
  \item the resulting collective stance and any changes to $(V, C, P)$ in the
        research integrity ledger.
\end{itemize}

\paragraph{Implementation.}

\begin{itemize}
  \item Integrate FRFP-aware evaluation forms into journal, conference,
        or internal review systems;
  \item encourage reviewers to explicitly comment on:
        \begin{itemize}
          \item the explicit/tacit structure of the work
                (what is well formalised vs.\ what relies on tacit expertise);
          \item any suspected safe-horizon violations
                (e.g.\ over-reliance on AI proofs or simulations);
          \item provenance and authorship (clarity about AI assistance).
        \end{itemize}
  \item Use these logs to refine safe-horizon and admissibility parameters
        for future projects.
\end{itemize}

\subsection{Metrics and Monitoring for Scientific FRFP}

To monitor long-horizon runs, an implementation can define simple, domain-tailored
metrics, understood as observables on runs and boundary events.

\subsubsection*{Example Metric Families}

\paragraph{Re-grounding frequency.}
For each project, measure:

\begin{itemize}
  \item how often primary sources (original papers, raw data) are re-read;
  \item how often core definitions are revisited and potentially refined;
  \item how often proofs are re-derived without AI assistance.
\end{itemize}

\paragraph{Boundary density.}
Measure, per unit length of the run (or per session):

\begin{itemize}
  \item number of boundary events;
  \item distribution of their $(V, C, P)$ statuses;
  \item rate of later revisions to earlier boundary decisions.
\end{itemize}

\paragraph{AI reliance profile.}
For a given artefact or project:

\begin{itemize}
  \item fraction of edits labeled as AI-authored in the hybrid authorship log;
  \item concentration of AI involvement by artefact type
        (statements vs.\ proofs vs.\ exposition);
  \item correlation between AI involvement and later revisions or corrections.
\end{itemize}

None of these metrics define correctness on their own.
They serve as \emph{signals} about degradation, safe-horizon stress, and
possible protocol failures that merit human attention.

\subsection{Illustrative Pattern: AI-Assisted Theory Development}

To make the above more concrete, consider a stylised pattern for FRFP-aligned
theoretical work.

\subsubsection*{Phase F0/F1: Exploration and Proto-Formalization}

\begin{itemize}
  \item Researchers use AI systems to survey literature, generate examples,
        propose conjectures, and suggest potential proof strategies.
  \item The run is tagged as exploratory (phases F0/F1), with relatively
        loose safe horizons and frequent AI involvement.
  \item There are few boundary events; results are treated as speculative.
\end{itemize}

\subsubsection*{Phase F2: Formal Development}

\begin{itemize}
  \item Selected conjectures are promoted to formal objects in $E$ and given
        stable labels.
  \item AI systems may still propose steps, but each non-trivial step in a
        proof must be locally checked by a human or an independent proof
        assistant; unchecked AI-generated steps are not admissible for boundary
        endorsement.
  \item Boundary events occur when the team accepts a lemma or theorem as
        provisionally correct.
        These events are logged with:
        the proof sketch, authorship information, and reviewer identity.
  \item Safe horizons constrain how many AI-suggested steps can be accepted
        without a fresh, human-driven re-derivation.
\end{itemize}

\subsubsection*{Phase F3: Consolidation and Dissemination}

\begin{itemize}
  \item Proofs are refactored into canonical forms, with minimised dependence
        on AI-generated scaffolding.
  \item Manuscripts are drafted with clear provenance for AI assistance.
  \item External peer review provides institutional boundary events, updating
        the research integrity ledger and potentially modifying $(V, C, P)$.
\end{itemize}

At every stage, FRFP's explicit/tacit split and boundary semantics guide tool
design (what AI is allowed to do), protocol design (where and how humans must
endorse), and governance (how committees and journals record and justify their
decisions).

\subsection{Summary}

This section has described an FRFP implementation tailored to scientific
research:

\begin{itemize}
\item explicit pipelines and artefact histories capture long-horizon research activity, including AI-assisted steps;
  \item tacit space $T$ and boundary events capture human scientific judgment
        and correctness;
  \item safe horizons and re-grounding practices control degradation and
        long-run hallucination risk;
  \item instruments (integrity ledgers, authorship logs, evaluation logs)
        and metrics make the structure of inquiry visible to researchers,
        reviewers, and institutions.
\end{itemize}

This implementation is intended as one concrete instantiation of FRFP in a
single domain (scientific inquiry).
Other domains---such as engineering, finance, or clinical practice---may
require different instruments and metrics, but can follow the same pattern:
start from the FRFP kernel, then design domain-specific pipelines, boundaries,
and governance on top.
% =========================================================
% Appendix: Implementing FRFP in Engineering Workflows
% =========================================================

\section{Implementing FRFP in Engineering Workflows}
\label{sec:frfp-engineering}

This section sketches a concrete implementation of the Foundational Reasoning
and Feedback Protocol (FRFP) for engineering domains, such as software,
infrastructure, and systems engineering.
We assume the FRFP core from the main text: an explicit/tacit split, long-horizon use with accumulating
effects, and a one-way handoff from explicit artefacts to human evaluation and endorsement. The goal here
is to make engineering runs auditable: what the system produced, what humans reviewed, and what was approved for action.

\subsection{Scope and Objectives}

We consider settings where:

\begin{itemize}
  \item AI systems assist in design, analysis, coding, verification,
        documentation, and incident analysis;
  \item correctness is grounded in specifications, physical constraints,
        safety standards, and organisational requirements;
  \item work is iterative and multi-stage: architecture, implementation,
        testing, deployment, and operations.
\end{itemize}

An FRFP-aligned engineering implementation should:

\begin{enumerate}
  \item keep AI confined to explicit artefacts (requirements drafts, designs,
        code, tests, logs);
  \item attach correctness only at human boundary events
        (approvals, sign-offs, risk acceptances);
  \item design safe horizons for automated transformations and refactors;
  \item surface provenance and authorship for later audit and incident review.
\end{enumerate}

\subsection{Core Principles for Engineering Pipelines}

\subsubsection*{E1: Explicit Engineering Artefacts}

Explicit space $E$ contains:

\begin{itemize}
  \item requirements documents, design specs, interface definitions;
  \item code, configuration files, infrastructure-as-code, models, and schemas;
  \item test cases, proofs, static-analysis and verification artefacts;
  \item logs, metrics, incident reports, and runbooks.
\end{itemize}

AI systems may:

\begin{itemize}
  \item generate and refactor code;
  \item propose designs and architectures;
  \item suggest test cases and checks;
  \item summarize logs and incident timelines.
\end{itemize}

All such contributions are explicit artefact transformations: drafts, diffs, refactors, test suggestions,
and summaries. They can be reviewed, checked, logged, and compared.

\subsubsection*{E2: Tacit Engineering Judgment and Boundary Events}

Tacit space $T$ contains:

\begin{itemize}
  \item engineers' understanding of the system, failure modes, and domain;
  \item awareness of regulatory, safety, and organisational constraints;
  \item practical sense of acceptable risk and robustness.
\end{itemize}

Boundary events in engineering include:

\begin{itemize}
  \item approval of system requirements or architecture;
  \item code review approvals and merge decisions;
  \item sign-off on test coverage and verification status;
  \item deployment approvals, change management sign-offs;
  \item acceptance of incident postmortems and remediation plans.
\end{itemize}

At each boundary, responsible humans in $T$ evaluate explicit artefacts in
$E$ and attach correctness for action.

\subsubsection*{E3: Safe Horizons in Automation and Refactoring}

Safe horizons for engineering pipelines limit:

\begin{itemize}
  \item how many automated refactors or transformations may be chained without
        human review;
  \item how far AI-assisted design exploration can go before re-grounding in
        core constraints and requirements;
  \item how long monitoring and remediation workflows can run under automated
        control before escalation.
\end{itemize}

An implementation must parameterise these horizons per system criticality
(e.g.\ safety-critical vs.\ internal tooling) and make the triggers explicit:
what forces escalation, rollback, or freeze of further automated changes.

\subsection{Engineering Methodology as FRFP Pipelines}

\subsubsection*{EM1: Requirements and Architecture}

\paragraph{Pipeline structure.}
A project begins with a run segment in which:

\begin{itemize}
  \item stakeholders articulate goals and constraints (explicit artefacts);
  \item AI systems may propose requirement drafts, refine wording, and check
        consistency with existing documents;
\item the workflow moves between stakeholder views and abstraction levels
        (user needs, non-functional requirements, threat model, interfaces).
\end{itemize}

\paragraph{Boundary events.}
Boundary events at this phase include:

\begin{itemize}
  \item formal approval of a requirements baseline;
  \item acceptance of an architecture document as the basis for implementation.
\end{itemize}

At these boundaries, humans must see:

\begin{itemize}
  \item the proposed document(s);
  \item provenance information indicating where AI contributed;
  \item trace links to standards, prior decisions, and risk assessments.
\end{itemize}

\paragraph{Implementation requirement.}
\begin{itemize}
  \item Treat requirements/architecture approval as an endorsement boundary that emits a durable record:
        approved version ID, approver role, known assumptions, and explicit “out of scope” exclusions.
\end{itemize}
\subsubsection*{EM2: Implementation and Verification}

\paragraph{Pipeline structure.}
The implementation phase is represented by a run segment
comprising:

\begin{itemize}
  \item explicit code changes (possibly AI-generated);
  \item test design and updates;
  \item static analysis, verification runs, and simulation artefacts.
\end{itemize}

\paragraph{Safe horizons.}
Safe horizons constrain:

\begin{itemize}
  \item length of AI-driven code modification chains between human reviews;
  \item number of test case generations before a human must inspect for
        coverage gaps and specification mismatch;
  \item amount of automatic config or IaC changes before manual audit.
\end{itemize}

\paragraph{Boundary events.}
Typical boundaries:

\begin{itemize}
  \item code review and merge approvals;
  \item verification and test sign-offs;
  \item internal releases and promotions between environments.
\end{itemize}

Each boundary is logged with:

\begin{itemize}
  \item the diff and tests considered;
  \item responsible reviewers and their judgments;
  \item any explicit delegation to downstream checks.
\end{itemize}

\subsubsection*{EM3: Deployment, Monitoring, and Incidents}

\paragraph{Pipeline structure.}
Deployment and operations produce runs over:

\begin{itemize}
  \item deployment artefacts and configuration states;
  \item monitoring dashboards, alerts, and traces;
  \item incident tickets, runbooks, and remediation changes.
\end{itemize}

AI may:

\begin{itemize}
  \item suggest rollouts and rollbacks;
  \item summarise incidents;
  \item propose remediation steps;
  \item cluster and interpret logs.
\end{itemize}

\paragraph{Boundary events.}
Boundary events include:

\begin{itemize}
  \item authorisation of production deployments and rollbacks;
  \item acceptance of remediation steps in incidents;
  \item closure of incidents and sign-off on postmortems.
\end{itemize}

These events must remain human decisions, potentially informed by AI
summaries but not replaced by them.

\subsection{Instruments and Metrics for Engineering FRFP}

\subsubsection*{EI1: Change and Review Ledger}

A change and review ledger records:

\begin{itemize}
  \item each change set (code/config/design) as an explicit transformation step;
  \item the associated tests and verification artefacts;
  \item review decisions and reviewers' identities;
  \item provenance of AI assistance in each change.
\end{itemize}

It enables queries such as:

\begin{itemize}
  \item ``Which production incidents involved changes with high AI authorship
        and minimal human review?''
  \item ``Where did we accept AI-suggested changes at shallow safe horizons?''
\end{itemize}

\subsubsection*{EI2: AI-Aided Change Profile}

Metrics include:

\begin{itemize}
  \item fraction of changes initiated by AI suggestions;
  \item distribution of AI involvement by subsystem or component;
  \item correlation between AI-involved changes and later defects or incidents.
\end{itemize}

These are not correctness guarantees, but signals for FRFP-aligned governance
to adjust safe horizons and review policies.

\subsection{Summary}

In engineering, FRFP yields an implementation in which:

\begin{itemize}
  \item AI operates on explicit artefacts (requirements, code, tests, logs);
  \item correctness is conferred only at human boundary events;
  \item safe horizons control automated change sequences and incident response;
  \item ledgers and logs make explicit pipelines and endorsements auditable.
\end{itemize}

% =========================================================
% Appendix: Implementing FRFP in Financial Decision-Making
% =========================================================

\section{Implementing FRFP in Financial Decision-Making}
\label{sec:frfp-finance}

This section outlines an implementation of FRFP for financial domains:
banking, trading, risk management, and related activities in which AI
systems provide analysis and recommendations while human agents and
institutions retain ultimate responsibility.

\subsection{Scope and Objectives}

We consider workflows where:

\begin{itemize}
  \item AI systems generate or transform credit analyses, risk reports,
        stress scenarios, trading ideas, and portfolio diagnostics;
  \item decisions have regulatory, fiduciary, or systemic implications;
  \item workflows are long-horizon (ongoing monitoring, repeated committee
        reviews, annual models and stress programmes).
\end{itemize}

An FRFP-compliant implementation should:

\begin{enumerate}
  \item confine AI to explicit artefacts (analyses, scenarios, reports,
        dashboards);
  \item reserve risk-taking and approvals as human boundary events;
  \item define safe horizons for AI-supported monitoring and strategy
        revisions;
  \item enable rigorous audit and regulatory scrutiny over explicit pipelines
        and decisions.
\end{enumerate}

\subsection{Core Principles for Financial Pipelines}

\subsubsection*{F1: Explicit Financial Artefacts}

Explicit space $E$ includes:

\begin{itemize}
  \item counterparty data, financial statements, market data, scenario sets;
  \item risk models, pricing models, stress-test engines, and their outputs;
  \item AI-generated narratives (credit memos, stress explanations, dashboards);
  \item policies, limits, risk appetite statements, and committee minutes.
\end{itemize}

AI may:

\begin{itemize}
  \item summarise exposures and drivers;
  \item construct scenario narratives and structured stress sets;
  \item generate draft memos and recommendation templates;
  \item flag anomalies and pattern shifts in monitoring data.
\end{itemize}

\subsubsection*{F2: Tacit Financial Judgment and Boundary Events}

Tacit space $T$ contains:

\begin{itemize}
  \item portfolio managers' and risk officers' understanding of markets,
        counterparties, products, and regulatory context;
  \item knowledge of past crises, liquidity dynamics, and behavioural effects;
  \item shared institutional norms about acceptable risk.
\end{itemize}

Boundary events include:

\begin{itemize}
  \item trade approvals and investment decisions;
  \item credit committee approvals and limit settings;
  \item model and scenario approvals by risk committees;
  \item acceptance of capital plans and regulatory submissions.
\end{itemize}

At each boundary, human decision-makers evaluate explicit artefacts and
attach correctness in light of $T$.

\subsubsection*{F3: Safe Horizons in Monitoring and Strategy}

Safe horizons specify:

\begin{itemize}
  \item how many AI-generated monitoring reports can be accepted without
        a full human re-baselining of assumptions and models;
  \item how often AI-suggested scenario updates require committee-level
        review rather than desk-level acceptance;
  \item what thresholds in metrics trigger escalated boundary events.
\end{itemize}

An implementation sets shorter horizons for high-risk or illiquid positions,
and longer ones for routine exposures.
\paragraph{Model governance alignment.}
\begin{itemize}
  \item Treat any AI system that affects risk-taking, limits, or regulatory outputs as a “model” for
        governance purposes: require documentation of intended use, limitations, validation evidence,
        ongoing monitoring, and escalation paths.
\end{itemize}
\subsection{Financial Methodology as FRFP Pipelines}

\subsubsection*{FM1: Credit and Investment Decisions}

\paragraph{Pipeline structure.}
For a credit or investment case, a run segment includes:

\begin{itemize}
  \item data gathering and validation;
  \item model runs (PD, LGD, PFE, etc.) and stress simulations;
  \item AI-generated summaries and draft memos;
  \item human edits, analysis, and commentary.
\end{itemize}

\paragraph{Boundary events.}
Boundary events include:

\begin{itemize}
  \item desk-level recommendation (accept/decline/hold);
  \item credit committee decision;
  \item limit-setting decisions for the counterparty or product.
\end{itemize}

Each event is logged with:

\begin{itemize}
  \item the artefacts considered (models, scenarios, memos);
  \item voting patterns or individual approvals;
  \item provenance of AI assistance in the written record.
\end{itemize}

\paragraph{Implementation requirement.}
\begin{itemize}
  \item Require an “assumption and sensitivity” block at endorsement boundaries (committee or desk lead):
        what drivers the decision depends on, what would change the decision, and which signals trigger
        re-review within the safe horizon.
\end{itemize}
\subsubsection*{FM2: Stress Testing and Capital Planning}

\paragraph{Pipeline structure.}
A stress-testing programme yields a run segment involving:

\begin{itemize}
  \item generation and refinement of macro and idiosyncratic scenarios;
  \item application of models and aggregations;
  \item AI-generated scenario narratives and impact summaries;
  \item human adjustments and overlays.
\end{itemize}

\paragraph{Safe horizons.}
Safe horizons limit:

\begin{itemize}
  \item the number of purely AI-generated scenario modifications accepted
        without independent macro or risk review;
  \item the depth of automated overrides of model outputs before escalation.
\end{itemize}

Boundary events occur when:

\begin{itemize}
  \item scenarios and methodologies are approved for use;
  \item final stress results are accepted for internal or regulatory use.
\end{itemize}

\subsubsection*{FM3: Ongoing Monitoring}

\paragraph{Pipeline structure.}
Monitoring yields repeated runs over:

\begin{itemize}
  \item time series of exposures and risk indicators;
  \item AI-generated anomaly flags, trend analyses, and alerts;
  \item human investigations and follow-ups.
\end{itemize}

\paragraph{Boundary events.}
Boundary events include:

\begin{itemize}
  \item escalation of a monitoring signal to a case or committee;
  \item changes in limits or hedging strategies;
  \item decisions to ignore or downgrade an AI-generated concern.
\end{itemize}
\paragraph{Escalation discipline.}
\begin{itemize}
  \item Make “ignore/downgrade” a logged boundary decision with rationale category (data quality issue,
        known model limitation, non-material exposure, false positive). This supports later audit and avoids
        silent normalization of alerts.
\end{itemize}
\subsection{Instruments and Metrics for Financial FRFP}

\subsubsection*{FI1: Decision and Scenario Ledger}

A decision ledger records:

\begin{itemize}
  \item the explicit pipeline leading to each credit or investment decision;
  \item the scenario and model set in force at that time;
  \item individual votes and reasoning;
  \item indicators of AI involvement in summaries and drafts.
\end{itemize}

\subsubsection*{FI2: Monitoring Profile}

Monitoring metrics include:

\begin{itemize}
  \item number and type of AI-generated alerts per period;
  \item fraction escalated to human review;
  \item fraction overridden by humans;
  \item time between repeated re-baselining of monitoring assumptions.
\end{itemize}

These metrics inform adjustments to safe horizons and escalation rules.

\subsection{Summary}

In finance, FRFP yields an implementation in which:

\begin{itemize}
  \item AI systems generate and process explicit artefacts, but do not
        autonomously take risk or approve positions;
  \item human decision-makers attach correctness at clearly defined boundaries;
  \item safe horizons govern monitoring, scenario use, and AI reliance;
  \item ledgers and logs make decisions and their AI involvement reconstructible
        for internal governance and external regulators.
\end{itemize}

% =========================================================
% Appendix: Implementing FRFP in Clinical Practice
% =========================================================

\section{Implementing FRFP in Clinical Practice}
\label{sec:frfp-clinical}

This section presents an implementation of FRFP in clinical workflows,
such as hospital care, outpatient medicine, and clinical decision support.
We again assume the FRFP core from the main text.

\subsection{Scope and Objectives}

We consider settings where:

\begin{itemize}
  \item AI systems summarise patient histories, suggest differential
        diagnoses, propose treatment plans, and assist with documentation;
  \item clinical decisions have direct impacts on patient health and legal
        liability;
  \item workflows involve multiple hand-offs and teams (nurses, physicians,
        specialists, pharmacists, administrators).
\end{itemize}

An FRFP-compliant clinical implementation should:

\begin{enumerate}
  \item confine AI to explicit artefacts (notes, summaries, suggestions);
  \item leave diagnosis and treatment decisions as human boundary acts;
  \item manage safe horizons in repeated, AI-assisted documentation and
        decision support;
  \item ensure that provenance and responsibility are clear in the record.
\end{enumerate}

\subsection{Core Principles for Clinical Pipelines}

\subsubsection*{C1: Explicit Clinical Artefacts}

Explicit space $E$ contains:

\begin{itemize}
  \item structured patient data (vitals, labs, imaging results, medications);
  \item free-text notes (H\&P, progress notes, discharge summaries);
  \item guidelines, order sets, protocols, and clinical pathways;
  \item AI-generated summaries, differentials, risk scores, and recommendations.
\end{itemize}

AI systems may:

\begin{itemize}
  \item draft notes from structured data and prior text;
  \item summarise long chart histories;
  \item suggest diagnoses or management options;
  \item rank guideline or order-set options given the current context.
\end{itemize}

\subsubsection*{C2: Tacit Clinical Expertise and Boundary Acts}

Tacit space $T$ contains:

\begin{itemize}
  \item clinician expertise, pattern recognition, and sense of atypicality;
  \item knowledge of local practice norms and patient-specific context;
  \item ethical judgment and responsibility for patient outcomes.
\end{itemize}

Boundary events include:

\begin{itemize}
  \item signing of notes and orders (diagnostic and therapeutic);
  \item finalising discharge summaries and hand-off documents;
  \item multidisciplinary conference decisions (tumour boards, complex case
        conferences);
  \item acceptance of AI-supported predictions into standard workflows.
\end{itemize}

At each boundary, a clinician (or a team) must evaluate AI contributions in
light of $T$ and attach correctness for action.
\paragraph{Transparency at the boundary.}
\begin{itemize}
  \item At endorsement points, the clinician must be able to see the basis of AI suggestions in a form that is
        clinically reviewable (source data referenced, guideline link, or stated rationale), and must be able to
        override. This keeps “decision support” from becoming de facto automation and aligns with common CDS governance
        expectations.
\end{itemize}
\subsubsection*{C3: Safe Horizons in Documentation and Decision Support}

Safe horizons control:

\begin{itemize}
  \item how many consecutive notes may be heavily AI-assisted without
        independent re-reading of primary data;
  \item how often AI-supported diagnostic suggestions must be audited against
        full guideline review or manual literature review;
  \item how long a patient can be managed with AI-summarised charts before
        a clinician must reconstruct key elements from raw data.
\end{itemize}

Horizon parameters depend on patient acuity, setting (ICU vs.\ outpatient),
and risk tolerance.

\subsection{Clinical Methodology as FRFP Pipelines}

\subsubsection*{CM1: Intake and Initial Assessment}

\paragraph{Pipeline structure.}
For a new admission, a run segment may include:

\begin{itemize}
  \item ingestion of prior records and structured data;
  \item AI-generated summary of history and active problems;
  \item clinician review, correction, and augmentation of the summary.
\end{itemize}

\paragraph{Boundary event.}
Boundary occurrence:

\begin{itemize}
  \item clinician signs the initial note and places initial orders.
\end{itemize}

At this point, correctness is conferred on:

\begin{itemize}
  \item the problem list and working diagnoses;
  \item key decisions (admission level, initial tests, immediate treatment).
\end{itemize}

\subsubsection*{CM2: Ongoing Management}

\paragraph{Pipeline structure.}
Over the course of care, a run segment includes:

\begin{itemize}
  \item repeated AI-assisted progress notes;
  \item AI-generated trend summaries (vitals, labs, risk scores);
  \item human interpretation, updates to diagnoses, and treatment adjustments.
\end{itemize}

\paragraph{Safe horizons.}
Safe horizons might enforce:

\begin{itemize}
  \item a limit on how many successive notes can be accepted with minimal
        editing of AI output;
  \item periodic requirements to review raw labs, images, and data streams
        directly;
  \item checks before relying on AI-generated risk scores for escalation
        decisions.
\end{itemize}

\paragraph{Boundary events.}
Boundary points include:

\begin{itemize}
  \item daily sign-off on progress notes and orders;
  \item decisions to escalate care or transfer units;
  \item significant changes in diagnosis or treatment plan.
\end{itemize}

\subsubsection*{CM3: Discharge and Handoffs}

\paragraph{Pipeline structure.}
At discharge or hand-off, a run segment contains:

\begin{itemize}
  \item AI-assisted draft of discharge summary or hand-off note;
  \item clinician edits to ensure accuracy, emphasising key risks and follow-up;
  \item generation of patient-facing summaries and instructions.
\end{itemize}

\paragraph{Boundary event.}
The final discharge or hand-off note is a boundary artefact:

\begin{itemize}
  \item its signing marks the clinician’s endorsement of its content;
  \item errors at this point can propagate to downstream providers and
        future care decisions.
\end{itemize}

FRFP requires that AI contributions are visible and that clinicians have
sufficient time and support to re-ground their tacit understanding before
signing.
\paragraph{Implementation requirement.}
\begin{itemize}
  \item Treat discharge and handoff sign-off as a protected boundary: the system must surface discrepancies (med lists,
        pending labs, contraindications, follow-up gaps) and must record whether clinician confirmation occurred.
\end{itemize}
\subsection{Instruments and Metrics for Clinical FRFP}

\subsubsection*{CI1: Clinical Documentation Ledger}

A documentation ledger records, for each patient:

\begin{itemize}
  \item all signed notes and orders as boundary artefacts;
  \item AI involvement in each note (e.g.\ percentage of AI-generated text,
        suggestions taken vs.\ rejected);
  \item episodes where clinicians explicitly overrode AI suggestions;
  \item links to underlying data and guidelines referenced.
\end{itemize}

This enables later review of:

\begin{itemize}
  \item how AI assistance correlates with documentation errors;
  \item patterns of overreliance or under-use across clinicians or units.
\end{itemize}

\subsubsection*{CI2: Decision-Support Profile}

For AI-supported decision aids (e.g.\ risk scores, triage suggestions),
metrics include:

\begin{itemize}
  \item rate of use and override by clinicians;
  \item cases where AI disagreed with clinician judgment and outcome data;
  \item time since last full audit of model performance and assumptions.
\end{itemize}

These metrics help calibrate safe horizons for decision support.

\subsection{Summary}

In clinical practice, FRFP yields an implementation in which:

\begin{itemize}
  \item AI systems operate on explicit artefacts (notes, summaries, scores);
  \item clinicians retain responsibility for diagnosis and treatment at
        boundary events;
  \item safe horizons limit extended, ungrounded reliance on AI outputs;
  \item documentation and evaluation logs make AI involvement and human
        judgment reconstructible for quality improvement and oversight.
\end{itemize}
